\manualchapter{ROMs and compatibility}{sec:roms}
\subsection{Loading a game}
Choose \textbf{Load ROM} from the module's right-click menu, or drag a ROM
file onto the panel. The chooser accepts \texttt{.nes} and \texttt{.NES}.
Extract compressed archives first.

Loading a ROM, including the same ROM again, resets the emulator and clears
SAVE. Canceling leaves the current selection alone. Only the first dropped
file is used.

\subsection{Supported cartridge mappers}
A mapper arranges cartridge memory. RackNES supports these six IDs;
individual games and modified ROMs may still behave incorrectly.
\begin{center}
\begin{tabular}{@{}rl@{\hspace{3em}}rl@{}}
\toprule
\textbf{ID} & \textbf{Name} & \textbf{ID} & \textbf{Name} \\
\midrule
0 & NROM & 3 & CNROM \\
1 & MMC1 & 7 & AxROM / AOROM \\
2 & UNROM & 9 & MMC2 / PxROM \\
\bottomrule
\end{tabular}
\end{center}
Mapper 7 accepts 32--256 KiB of PRG ROM in power-of-two sizes, 8 KiB of
CHR RAM, and either one-screen nametable page. Legacy iNES and NES 2.0
submappers 0/1 use no bus conflicts; submapper 2 combines writes with the
visible PRG byte. Only NTSC layouts are accepted. CHR ROM, battery RAM,
trainers, four-screen layouts, extra devices, and other variants are rejected.
Legacy headers cannot reliably identify every board variant.

Mapper 9 accepts power-of-two PRG ROM sizes of 32--128 KiB, CHR ROM sizes
of 8--128 KiB, and 8 KiB PRG RAM, with optional battery metadata. Supported
formats are NTSC iNES and NES 2.0 submapper 0. Trainers, CHR RAM, four-screen
layouts and other variants are rejected. RAM persists in Rack patches and
SAVE snapshots; external battery-save files are not imported or exported.

Timing is fixed; there is no PAL/NTSC or mapper selector. Audio uses the five
base NES voices without cartridge expansion sound. CV Genie's game list
contains memory maps and does not define ROM compatibility.

\subsection{When loading fails}
\begin{description}
\item[ASIC mapper not implemented for ROM!]
The file requests a mapper outside the implemented set, or an unsupported
mapper-7 or mapper-9 image layout.
\item[ROM file failed to load!]
The selected path failed ROM validation. Check the file and its format.
\item[ROM file was not found!]
A patch's stored ROM path could not be validated on restore. Check that the
same ROM is still present at that path.
\end{description}
A failed new-ROM attempt leaves the current game, display, and SAVE slot
in place. Correct the file selection and try again.

\textbf{Moving patches.} Keep the original ROM at its saved path: restoration
reloads that file. Loading a ROM manually starts a new session; it does not
recover a failed snapshot restore.
